% GENERATED FILE. Edit canonical lessons, then run npm run generate:books.
% canonical-source-hash: fnv1a64:2362fd508d7a6c09
% canonical-lessons: KA-C208-deha, KA-C208-cuccumaddu, KA-C208-mulamu, KA-C208-kayile, KA-C208-rogi

\chapter{Body, Injection, Ointment, Illness, Patient}
\label{ch:ka-body-injection-ointment-illness-patient}
\hlchaptermodality{208}

\begin{chapteropening}
\textbf{By the end of this chapter:} \emph{I can say the body, an injection, an ointment, an illness and a patient.}

Complete the last lesson of chapter 208: I can say the body, an injection, an ointment, an illness and a patient.
\end{chapteropening}

\section[\texorpdfstring{\kn{ದೇಹ} (dēha)}{ದೇಹ (dēha)}]{\kn{ದೇಹ} (dēha) — the body}

\label{lesson:KA-C208-deha}



\begin{quote}
Before the new one: say the Kannada for to stack, then the Kannada for to collect.
\end{quote}

\subsection*{You'll want to know: \kn{ದೇಹ}}
\textbf{\kn{ದೇಹ}} — \emph{dēha} — \textquotedblleft{}the body\textquotedblright{}.

\begin{grammarlens}[title={What you've built}]
One more word for this chapter.
\end{grammarlens}

\subsection*{Your turn}
\begin{itemize}
\raggedright
  \item \emph{Say it:} \emph{dēha}
  \item \emph{Say it:} \emph{dēha}, once more
  \item \emph{Say it:} say \emph{dēha}
  \item \emph{Recall:} say the Kannada for to stack, then the Kannada for to collect, then say \emph{dēha} again
\end{itemize}

\begin{tcolorbox}[breakable,colback=teal!4,colframe=teal!35!black,title={Before you move on}]
What does \kn{ದೇಹ} mean? (\textquotedblleft{}The body\textquotedblright{}.) Say it once more.
\end{tcolorbox}
\section[\texorpdfstring{\kn{ಚುಚ್ಚುಮದ್ದು} (cuccumaddu)}{ಚುಚ್ಚುಮದ್ದು (cuccumaddu)}]{\kn{ಚುಚ್ಚುಮದ್ದು} (cuccumaddu) — an injection}

\label{lesson:KA-C208-cuccumaddu}



\begin{quote}
Before the new one: say the Kannada for to collect, then the Kannada for the body.
\end{quote}

\subsection*{You'll want to know: \kn{ಚುಚ್ಚುಮದ್ದು}}
\textbf{\kn{ಚುಚ್ಚುಮದ್ದು}} — \emph{cuccumaddu} — \textquotedblleft{}an injection\textquotedblright{}.

\begin{grammarlens}[title={What you've built}]
One more word for this chapter.
\end{grammarlens}

\subsection*{Your turn}
\begin{itemize}
\raggedright
  \item \emph{Say it:} \emph{cuccumaddu}
  \item \emph{Say it:} \emph{cuccumaddu}, once more
  \item \emph{Say it:} say \emph{cuccumaddu}
  \item \emph{Recall:} say the Kannada for to collect, then the Kannada for the body, then say \emph{cuccumaddu} again
\end{itemize}

\begin{tcolorbox}[breakable,colback=teal!4,colframe=teal!35!black,title={Before you move on}]
What does \kn{ಚುಚ್ಚುಮದ್ದು} mean? (\textquotedblleft{}An injection\textquotedblright{}.) Say it once more.
\end{tcolorbox}
\section[\texorpdfstring{\kn{ಮುಲಾಮು} (mulāmu)}{ಮುಲಾಮು (mulāmu)}]{\kn{ಮುಲಾಮು} (mulāmu) — an ointment}

\label{lesson:KA-C208-mulamu}



\begin{quote}
Before the new one: say the Kannada for the body, then the Kannada for an injection.
\end{quote}

\subsection*{You'll want to know: \kn{ಮುಲಾಮು}}
\textbf{\kn{ಮುಲಾಮು}} — \emph{mulāmu} — \textquotedblleft{}an ointment\textquotedblright{}.

\begin{grammarlens}[title={What you've built}]
One more word for this chapter.
\end{grammarlens}

\subsection*{Your turn}
\begin{itemize}
\raggedright
  \item \emph{Say it:} \emph{mulāmu}
  \item \emph{Say it:} \emph{mulāmu}, once more
  \item \emph{Say it:} say \emph{mulāmu}
  \item \emph{Recall:} say the Kannada for the body, then the Kannada for an injection, then say \emph{mulāmu} again
\end{itemize}

\begin{tcolorbox}[breakable,colback=teal!4,colframe=teal!35!black,title={Before you move on}]
What does \kn{ಮುಲಾಮು} mean? (\textquotedblleft{}An ointment\textquotedblright{}.) Say it once more.
\end{tcolorbox}
\section[\texorpdfstring{\kn{ಕಾಯಿಲೆ} (kāyile)}{ಕಾಯಿಲೆ (kāyile)}]{\kn{ಕಾಯಿಲೆ} (kāyile) — an illness, a disease}

\label{lesson:KA-C208-kayile}



\begin{quote}
Before the new one: say the Kannada for an injection, then the Kannada for an ointment.
\end{quote}

\subsection*{You'll want to know: \kn{ಕಾಯಿಲೆ}}
\textbf{\kn{ಕಾಯಿಲೆ}} — \emph{kāyile} — \textquotedblleft{}an illness, a disease\textquotedblright{}.

\begin{grammarlens}[title={What you've built}]
One more word for this chapter.
\end{grammarlens}

\subsection*{Your turn}
\begin{itemize}
\raggedright
  \item \emph{Say it:} \emph{kāyile}
  \item \emph{Say it:} \emph{kāyile}, once more
  \item \emph{Say it:} say \emph{kāyile}
  \item \emph{Recall:} say the Kannada for an injection, then the Kannada for an ointment, then say \emph{kāyile} again
\end{itemize}

\begin{tcolorbox}[breakable,colback=teal!4,colframe=teal!35!black,title={Before you move on}]
What does \kn{ಕಾಯಿಲೆ} mean? (\textquotedblleft{}An illness, a disease\textquotedblright{}.) Say it once more.
\end{tcolorbox}
\section[\texorpdfstring{\kn{ರೋಗಿ} (rōgi)}{ರೋಗಿ (rōgi)}]{\kn{ರೋಗಿ} (rōgi) — a patient}

\label{lesson:KA-C208-rogi}



\begin{quote}
Before the new one: say the Kannada for an ointment, then the Kannada for an illness.
\end{quote}

\subsection*{You'll want to know: \kn{ರೋಗಿ}}
\textbf{\kn{ರೋಗಿ}} — \emph{rōgi} — \textquotedblleft{}a patient\textquotedblright{}.

\begin{grammarlens}[title={What you've built}]
One more word for this chapter.
\end{grammarlens}

\subsection*{Your turn}
\begin{itemize}
\raggedright
  \item \emph{Say it:} \emph{rōgi}
  \item \emph{Say it:} \emph{rōgi}, once more
  \item \emph{Say it:} say \emph{rōgi}
  \item \emph{Recall:} say the Kannada for an ointment, then the Kannada for an illness, then say \emph{rōgi} again
\end{itemize}

\begin{tcolorbox}[breakable,colback=teal!4,colframe=teal!35!black,title={Before you move on}]
What does \kn{ರೋಗಿ} mean? (\textquotedblleft{}A patient\textquotedblright{}.) Say it once more.
\end{tcolorbox}
